% Fuentes: EXP-005, EXP-013, EXP-014, EXP-015, EXP-017, EXP-018 ("Siguiente paso").
\section{Conclusiones y siguientes pasos}
\begin{itemize}
  \item \textbf{Lo que más aportó:} reducir la entrada a \res{final-res}~$\times$~\res{final-res}
  (+\res{gain-resolution}~pp), tratar el desbalance (+\res{gain-imbalance}~pp) y la invarianza a la orientación
  mediante augmentation D4 (+\res{gain-d4}~pp). En el modelo final, la profundidad y BatchNorm son además
  imprescindibles (Cuadro~\ref{tab:ablation}).
  \item \textbf{Lo que no aportó:} dropout, L2, schedules de la \emph{learning rate} y la mayoría de ajustes del
  optimizador, cuyas diferencias quedaron dentro del ruido entre semillas.
  \item \textbf{Límite:} una red densa sobre píxeles no comparte pesos entre posiciones ni tiene invarianzas
  espaciales; los vehículos, que a esta resolución se distinguen por detalles finos, siguen confundiéndose.
  \item \textbf{Siguiente paso:} redes convolucionales, que incorporan por diseño la conectividad local y la
  equivarianza a traslaciones~\cite{goodfellow}, partiendo de este modelo como baseline de la siguiente entrega.
\end{itemize}

Todas las cifras proceden de \texttt{experiments/experiments.csv} (\res{n-experiments} entrenamientos); el
razonamiento de cada decisión está en la bitácora del proyecto.

\begin{thebibliography}{9}\small
\bibitem{xview} D. Lam et al. \emph{xView: Objects in Context in Overhead Imagery}. arXiv:1802.07856, 2018.
\bibitem{lecun} Y. LeCun, L. Bottou, G. Orr, K.-R. Müller. \emph{Efficient BackProp}. En \emph{Neural Networks: Tricks of the Trade}, Springer, 1998.
\bibitem{goodfellow} I. Goodfellow, Y. Bengio, A. Courville. \emph{Deep Learning}. MIT Press, 2016.
\bibitem{ioffe} S. Ioffe, C. Szegedy. \emph{Batch Normalization: Accelerating Deep Network Training by Reducing Internal Covariate Shift}. ICML, 2015.
\bibitem{bouthillier} X. Bouthillier et al. \emph{Accounting for Variance in Machine Learning Benchmarks}. MLSys, 2021.
\bibitem{cawley} G. Cawley, N. Talbot. \emph{On Over-fitting in Model Selection and Subsequent Selection Bias in Performance Evaluation}. JMLR 11, 2010.
\end{thebibliography}
